%% ma: L11-B16-0080
%% lop: 11
%% chuong: 5
%% bai: 16
%% dang: 11.16.2
%% loai: TL
%% muc: TH
%% dap_an: "-"
%% nguon: "Ảnh giáo viên gửi 25/11/2026: Giới hạn hàm số dạng 0/0 (3 bài, có ĐS), Bài 1 câu 8"
%% kien_thuc: "Lớp 11 Bài 16 (giới hạn của hàm số tại một điểm); phân tích biểu thức chứa căn thành nhân tử"
%% trang_thai: chinh-thuc
%% ghi_chu: "Đề gốc ghi tử x - 2√x - 3: khi đó tử dần về -4, mẫu dần về 0, giới hạn không hữu hạn và không ra ĐS -4/3. Đã sửa tử thành x + 2√x - 3 để khớp ĐS."
\begin{ex}
Tính giới hạn $A=\lim\limits_{x\to1}\dfrac{x+2\sqrt x-3}{x-5\sqrt x+4}$.
\loigiai{$\dfrac{(\sqrt x-1)(\sqrt x+3)}{(\sqrt x-1)(\sqrt x-4)}=\dfrac{\sqrt x+3}{\sqrt x-4}$. Vậy $A=-\dfrac43$.}
\end{ex}
